% -------------------------------
% Plantilla del Anteproyecto en la Sede Regional Brunca. Versión 0.0 (Preliminar)
%
% CTFG Sede Regional Brunca
% 
% Compilar con XeLaTeX 
% -------------------------------


% -------------------------------
% Idioma
% -------------------------------
\usepackage{polyglossia}
\setmainlanguage{spanish}


% -------------------------------
% Fuente
% -------------------------------
% Cualquier tamaño de texto: {\fontsize{100pt}{120pt}\selectfont tutexto}
\usepackage{anyfontsize}
% Selección de fuentes.
\usepackage{fontspec}
% Espacio entre líneas
\usepackage{setspace}
% Fuentes especiales
\usepackage{textcomp,marvosym,pifont}


% -------------------------------
% Otros paquetes de uso común
% -------------------------------
% Símbolos matemáticos
\usepackage{amsmath,amsfonts,amssymb}
% Gráficos
\usepackage{graphicx}
% Posición arbitraria de los gráficos
\usepackage[absolute,overlay]{textpos}
% Apéndices
\usepackage{appendix}


% -------------------------------
% Esquema de colores
% -------------------------------
% Paquete para la definición de colores
\usepackage[table]{xcolor}
% Esquema de colores
\input{include/colores}


% -------------------------------
% Márgenes
% -------------------------------
% Márgenes de las páginas
\usepackage[
	inner	 =	3cm,   % Margen interior
	outer	 =	2.5cm, % Margen exterior
	top	 =	2.5cm, % Margen superior
	bottom =	2.5cm, % Margen inferior
	includeheadfoot, % Incluye cabecera y pie de página en los márgenes
]{geometry}
% Muestra una regla para comprobar el formato de las páginas (descomentar)
%\usepackage[type=upperleft,showframe,marklength=8mm]{fgruler} 


% -------------------------------
% Renombra tablas y bibliografía
% -------------------------------
\addto\captionsspanish{
	\renewcommand{\listtablename}{Índice de tablas}
	\renewcommand{\tablename}{Tabla}
	\renewcommand{\bibname}{Referencias bibliográficas}
}


% -------------------------------
% Itemize y enumerate. Control de los espacios.
% -------------------------------
\usepackage{enumitem}
\setitemize{itemsep=0pt}
\setenumerate{itemsep=0pt}
\renewcommand{\labelitemi}{\tiny\bft{$\blacksquare$}}
\renewcommand{\labelitemii}{\bft{$\bullet$}}
\renewcommand{\labelitemiii}{\bft{-}}


% -------------------------------
% Enlaces y colores
% -------------------------------
\usepackage{hyperref}
%\hypersetup{
%    colorlinks,
%    citecolor=black,
%    filecolor=black,
%    linkcolor=black,
%    urlcolor=black
%}

% -------------------------------
% Gestión de elementos flotantes
% -------------------------------
% Para forzar a que las figuras aparezcan antes del fin de una sección.
%\usepackage[section]{placeins}

% Gestión de títulos de los elementos flotantes
\usepackage{caption}
% Subfiguras y títulos en subfiguras
\usepackage{subcaption}

% Títulos
\DeclareCaptionFont{tema}{\color{tema}} % Color
\captionsetup{labelfont={tema, bf}}     % Estilo
\captionsetup{font=normalsize}          % Tamaño
\captionsetup{width=.9\linewidth}       % Anchura del título
% Distancia de la figura al texto.
\setlength{\intextsep}{1cm}
\setlength{\textfloatsep}{1cm}


% -------------------------------
% Utilidades para tablas
% -------------------------------
% Para fundir filas en tablas
\usepackage{multirow}
% Para definir columnas con ancho fijo y alineación
\usepackage{array,ragged2e}
\newcolumntype{L}[1]{>{\raggedright\let\newline\\\arraybackslash\hspace{0pt}}m{#1}}
\newcolumntype{C}[1]{>{\centering\let\newline\\\arraybackslash\hspace{0pt}}m{#1}}
\newcolumntype{R}[1]{>{\raggedleft\let\newline\\\arraybackslash\hspace{0pt}}m{#1}}


% -------------------------------
% Para tabular
% -------------------------------
\usepackage{tabto}


% -------------------------------
% Formato de títulos, capítulos, etc.
% -------------------------------
\usepackage{titlesec, titletoc}
% Parte
\titleformat{\part}[display]												   	   % Estilo
{\Huge \centering \color{tema}}                                                % Formato
{Parte \thepart }                                                              % Etiqueta
{0.75ex}                                                                       % Separación
{\bfseries\fontsize{25pt}{40pt}\selectfont}                                    % Entre etiqueta y título            
[\thispagestyle{empty}]                                                        % Después

% Capítulo		
\titleformat{\chapter}[block]            			 			                    % Estilo  	
{\vspace{0cm} \flushright \bfseries \LARGE \color{tema}}  	        			% Formato 
{\thechapter.}                  			  									% Etiqueta
{0.75ex}                                     									% Separación
{}                                           									% Entre etiqueta y título
[\vspace{0.1cm} \rule{0.5\textwidth}{1pt}\vspace*{1cm}\thispagestyle{plain}]    % Después      

% Sección
\titleformat{\section}
{\vspace{5pt}  \Large \color{tema}}
{\thesection .}
{0.75ex}
{}

% Subsección
\titleformat{\subsection}
{\large \color{tema}}
{\thesubsection .}
{0.75ex}
{}

% Subsubsección
\titleformat{\subsubsection}
{\bf\color{tema}}
{}
{0.75ex}
{}

% Parte (TOC)
\titlecontents{part}[2pc]
{\vspace{20pt}\bfseries\filright\Large\color{tema}}
{\contentslabel{1.5pc}}{\hspace*{-1.5pc}}
{\vspace{5pt}}
{}

% Capítulo (TOC)
\titlecontents{chapter}[2pc]
{\vspace{10pt}\bfseries\large\filright}
{\contentslabel{1.5pc}}{\hspace*{-1.5pc}}
{\mdseries\titlerule*[0.5pc]{.}\bfseries\contentspage}

% Sección (TOC)
\titlecontents{section}[4pc]
{\vspace{5pt}\filright}
{\contentslabel{2pc}}{\hspace*{2pc}}
{\titlerule*[0.5pc]{.}\contentspage}
{\vspace{5pt}

	% Subsección (TOC)
	\titlecontents{subsection}[6pc]
	{\vspace{2.5pt}\filright\small\em}
	{\contentslabel{3pc}}{\hspace*{-3pc}}
	{\titlerule*[0.5pc]{.}\contentspage}}


% -------------------------------
% Cabeceras y pie de página
% -------------------------------
\usepackage{fancyhdr}

\pagestyle{fancy}
\fancyhf{} % Limpiar todas las cabeceras y pies de página

% Definir el nombre de la empresa
\newcommand{\nombreempresa}{Café y Vivero Verde Pradera}

% Pie de página centrado con el nombre de la empresa
\fancyfoot[C]{\nombreempresa}

% Cabecera con el nombre del capítulo a la izquierda
\fancyhead[L]{\leftmark}  % Título del capítulo
\fancyhead[R]{\thepage}   % Número de página

% Ajuste de las líneas de cabecera y pie de página
\renewcommand{\headrulewidth}{0.5pt}  % Línea en la cabecera
\renewcommand{\footrulewidth}{0.5pt}  % Línea en el pie de página

% Redefinición del estilo de página 'plain' para capítulos y partes
\fancypagestyle{plain}{
	\fancyhf{}  % Limpiar cabeceras y pies de página
	\fancyfoot[LE,RO]{\thepage} % Número de página a la izquierda en páginas pares, derecha en impares
	\fancyfoot[C]{\nombreempresa}  % Nombre de la empresa en el pie de página
	\renewcommand{\headrulewidth}{0pt}  % Sin línea en la cabecera
	\renewcommand{\footrulewidth}{0.5pt}  % Línea en el pie de página
}

% -------------------------------
% Utilidades
% -------------------------------
% Marcas
\newcommand{\pcite}{\bfr{[?]}\,} % Cita pendiente: [?] en rojo y negrita.
\newcommand{\ps}{\bfr{[-}\,} % Pricipio de un bloque en sucio
\newcommand{\fs}{\bfr{-]}\,} % Fin de un bloque en sucio.

% Notas
\usepackage[textsize=tiny,spanish,shadow,textwidth=2cm]{todonotes}

% Para generar textos "random"
\usepackage{lipsum}




% -------------------------------
% Algoritmos y código
% -------------------------------

% Se definen entornos específicos para algoritmos y código que permiten
% gestionar los títulos manera similar en ambos casos, y similar al resto
% de elementos flotantes.
\usepackage{newfloat}

% Define un estilo para el encabezado de algoritmos y código.
\DeclareCaptionFormat{algcode}{
	\rule{\dimexpr\textwidth\relax}{1pt}\vspace{0.2cm}
	#1#2#3
	\rule{\dimexpr\textwidth\relax}{1pt}
}
\captionsetup[algorithm]{format=algcode, width=\linewidth}
\captionsetup[code]{format=algcode, width=\linewidth}


% -------------------------------
% Caracteres pifont de uso común (con color)
% -------------------------------
\newcommand{\vmarkt}{\textcolor{tema}{\ding{52}} }  % Checked (V) Tema
\newcommand{\vmarkr}{\textcolor{rojo}{\ding{52}} }  % Checked (V) Rojo
\newcommand{\vmarka}{\textcolor{azul}{\ding{52}} }  % Checked (V) Azul

\newcommand{\xmarkt}{\textcolor{tema}{\ding{55}} }  % Checked (X) Tema
\newcommand{\xmarkr}{\textcolor{rojo}{\ding{55}} }  % Checked (X) Rojo
\newcommand{\xmarka}{\textcolor{azul}{\ding{55}} }  % Checked (X) Azul

\newcommand{\lmarkt}{\textcolor{tema}{\ding{46}} }  % Lápiz Tema
\newcommand{\lmarkr}{\textcolor{rojo}{\ding{46}} }  % Lápiz Rojo
\newcommand{\lmarka}{\textcolor{azul}{\ding{46}} }  % Lápiz Azul

\newcommand{\hmarkt}{\textcolor{tema}{\ding{43}} }  % Mano Tema
\newcommand{\hmarkr}{\textcolor{rojo}{\ding{43}} }  % Mano Roja 
\newcommand{\hmarka}{\textcolor{azul}{\ding{43}} }  % Mano Azula 

\newcommand{\fmarkt}{\textcolor{tema}{\ding{220}} }  % Flechza Tema
\newcommand{\fmarkr}{\textcolor{rojo}{\ding{220}} }  % Flecha Roja 
\newcommand{\fmarka}{\textcolor{azul}{\ding{220}} }  % Flecha Azul